% from ega4-13-fr.tex
\oldpage[IV-3]{192}
La comparaison de ces inégalités montre donc que~:

\begin{lemma}[13.2.5]
\label{IV.13.2.5-fr}
Avec les notations de (13.2.4), les conditions suivantes sont équivalentes~:
\begin{enumerate}[label=\emph{\alph*)}]
  \item $\dim(B')=\dim(A)+\dim(B\otimes_Ak(y))$.
  \item On a simultanément les relations suivantes~:
  \begin{enumerate}[label=(\roman*)]
    \item $\dim(A')=\dim(A)$.
    \item $\dim(B'\otimes_{A'}k(y))=\dim(B\otimes_Ak(y))$, autrement dit
    \[
      \dim_x(X'\cap f^{-1}(y))=\dim_x(f^{-1}(y)).
    \]
    \item $\dim_x(f'^{-1}(y))=\dim(f'^{-1}(\eta'))$, autrement dit
    $f':X'\to Y'$ est équidimensionnel (13.2.2) au point $x$.
    \item On a l'égalité
    \[
      \dim(B')=\dim(A')+\dim(B'\otimes_{A'}k(y))
      -(\dim_x(f'^{-1}(y))-\dim(f'^{-1}(\eta')))
    \]
  \end{enumerate}
  (relation qui est toujours vérifiée lorsque $A=\mathcal O_{Y,y}$ est un
  anneau universellement caténaire, en vertu de (5.6.5)).
  \item On a simultanément les relations suivantes~:
  \begin{enumerate}[label=(\roman*)]
    \item $\dim(B')=\dim(B)$.
    \item $\dim(B)=\dim(A)+\dim(B\otimes_Ak(y))$.
  \end{enumerate}
\end{enumerate}
\end{lemma}

\begin{env}[13.2.6]
\label{IV.13.2.6-fr}
Rappelons maintenant que les composantes irréductibles $X_i$ de $X$ contenant
$x$ sont en nombre fini et que l'on a ((5.1.2.1) et ($\mathbf 0$, 14.2.1.1))
\[
\tag{13.2.6.1}
\label{IV.13.2.6.1-fr}
  \dim(\mathcal O_{X,x})=\sup_i\dim(\mathcal O_{X_i,x}).
\]
L'équivalence des conditions $b)$ et $c)$ dans (13.2.5) implique par suite,
compte tenu de ($\mathbf 0$, 16.3.9) et ($\mathbf 0$, 14.2.1)~:
\end{env}

\begin{proposition}[13.2.7]
\label{IV.13.2.7-fr}
Soient $Y$ un préschéma localement noethérien, $f:X\to Y$ un morphisme de type
fini, $x$ un point de $X$, $f(x)=y$~; on a
\[
\tag{13.2.7.1}
\label{IV.13.2.7.1-fr}
  \dim(\mathcal O_{X,x})\leq\dim(\mathcal O_{Y,y})
  +\dim(\mathcal O_{X,x}\otimes_{\mathcal O_{Y,y}}k(y)).
\]
Pour que les deux membres de (13.2.7.1) soient égaux, il faut et il suffit
qu'il existe une partie fermée irréductible $X'$ de $X$ contenant $x$ et
vérifiant simultanément les conditions suivantes~:
\begin{enumerate}[label=(\roman*)]
  \item Si $Y'=\overline{f(X')}$, on a
  $\dim(\mathcal O_{Y',y})=\dim(\mathcal O_{Y,y})$.
  \item $\dim_x(X'\cap f^{-1}(y))=\dim_x(f^{-1}(y))$ (autrement dit, $X'$
  contient une des composantes irréductibles de $f^{-1}(y)$ contenant $x$, de
  dimension maxima parmi toutes ces composantes). Il revient au même de dire
  que
  $\dim(\mathcal O_{X',x}\otimes_{\mathcal O_{Y',y}}k(y))
  =\dim(\mathcal O_{X,x}\otimes_{\mathcal O_{Y,y}}k(y))$.
  \item $X'$ est équidimensionnel sur $Y'$ au point $x$, autrement dit, on a
  \[
    \dim_x(X'\cap f^{-1}(y))=\dim(X'\cap f^{-1}(\eta')),
  \]
  où $\eta'$ est le point générique de $Y'$ (et par suite, toutes les
  composantes irréductibles de $X'\cap f^{-1}(y)$ contenant $x$ ont même
  dimension (13.2.2)).

  